\documentclass[12pt,a4paper]{article}
\usepackage[UTF8]{ctex}
\usepackage{amsmath,amssymb,booktabs,graphicx,float,geometry,array,longtable}
\geometry{left=2.2cm,right=2.2cm,top=2.5cm,bottom=2.5cm}
\title{NIPT 的时点选择与胎儿异常判定研究}
\date{}

\begin{document}
\maketitle

\begin{abstract}
针对无创产前检测（NIPT）中检测时点选择与胎儿异常判定问题，本文基于 267 名男胎孕妇 1082 条检测记录和 147 名女胎孕妇 605 条检测记录，建立了 Y 染色体浓度关系模型、分组时点优化模型以及女胎异常判定模型。首先，对男胎样本进行相关性分析与回归建模，发现 Y 染色体浓度与孕周呈显著正相关，与孕妇 BMI 呈显著负相关。两因素线性模型表明：Y 染色体浓度可表示为孕周的增函数和 BMI 的减函数，模型拟合结果中孕周项和 BMI 项均具有极高显著性。其次，以“Y 染色体浓度首次达到或超过 4\%”定义最早达标时间，构造基于 BMI 分组的时点优化模型，并以尽早发现异常和保证检测准确性之间的平衡为目标，对男胎孕妇进行合理分组。结果表明，当 BMI 采用 $[0,30)$、$[30,34)$、$[34,38)$ 和 $[38,+\infty)$ 的四组划分时，各组推荐检测时点分别可取约 14 周、14 周、17 周和 19 周附近。再次，在综合分析身高、体重、年龄等因素后发现，高 BMI 与高体重显著延迟达标时间，而年龄对 Y 浓度具有弱负向影响，因此第三问中在 BMI 分组基础上引入多因素修正可进一步提高时点决策的稳定性。最后，对于女胎异常判定问题，本文以 13、18、21 号染色体非整倍体标签为结果变量，分析 X 染色体浓度、Z 值及测序质量指标的判别作用，提出“X 染色体浓度阈值 + 染色体 Z 值阈值”的联合判别规则。结果显示，相较单纯使用 Z 值阈值，联合规则在召回率上有明显提升，更适用于国赛场景下的初筛判别。本文方法具有较强可解释性，可为不同 BMI 孕妇的 NIPT 时点制定和女胎异常判定提供量化依据。

\noindent\textbf{关键词：} NIPT；Y 染色体浓度；BMI 分组；检测时点优化；异常判定
\end{abstract}

\section{问题重述}
NIPT 是通过采集孕妇外周血、分析胎儿游离 DNA 片段，判断胎儿是否存在染色体异常的重要产前筛查技术。对于男胎而言，Y 染色体浓度是判断检测是否达到基本准确性要求的重要指标；若 Y 染色体浓度达到或超过 4\%，则通常认为检测结果较为可靠。对于女胎，由于不存在 Y 染色体，异常判定主要依赖 X 染色体及 13、18、21 号染色体相关指标。题目要求在给定检测数据基础上解决以下四个问题：

\begin{enumerate}
    \item 分析男胎 Y 染色体浓度与孕周、BMI 等指标之间的关系，并建立显著性模型；
    \item 仅依据 BMI 对男胎孕妇进行合理分组，给出各组最佳 NIPT 时点，使孕妇潜在风险最小；
    \item 综合身高、体重、年龄等因素，在考虑检测误差及达标比例的条件下，进一步优化男胎孕妇分组与最佳检测时点；
    \item 针对女胎数据，综合染色体 Z 值、GC 含量、读段数及其比例、BMI 等因素，建立异常判定方法。
\end{enumerate}

\section{问题分析}
本题的核心在于同时兼顾“尽早检测”和“检测可靠性”两种目标。若检测过早，则胎儿游离 DNA 信号不足，可能导致测序失败或 Y 染色体浓度达不到阈值；若检测过晚，虽然结果更稳定，但异常发现延迟会压缩治疗窗口期。因此，问题 2 和问题 3 本质上都是在“早检风险”和“漏检风险”之间求平衡。

从数据结构看，男胎和女胎样本均包含多次检测记录，部分孕妇存在“多次采血多次检测”或“一次采血多次检测”的情况。这意味着同一孕妇的多条记录并非独立样本，因此不能简单把全部记录等权视为不同个体。基于题意，我们将“最早达标时间”定义为同一孕妇所有检测记录中 Y 染色体浓度第一次达到 4\% 的最小孕周，这一处理既符合临床意义，也能避免重复检测对时点分析的干扰。

对于问题 1，可采用描述统计、相关性分析与回归模型建立 Y 浓度与各变量的关系。对于问题 2，可先按 BMI 候选区间分组，计算各组孕妇的最早达标时间分布，再根据较高达标率条件下的最早点作为推荐检测时点。对于问题 3，可在问题 2 的基础上，进一步把年龄、身高、体重等变量纳入修正项，从而实现更细粒度的时点决策。对于问题 4，由于女胎异常样本相对较少，应优先构造可解释、稳健的判别规则，而非依赖黑箱模型。

\section{模型假设}
\begin{enumerate}
    \item 同一孕妇不同检测记录反映的是同一妊娠过程中的动态变化，可用于估计最早达标时间；
    \item 将 Y 染色体浓度达到或超过 4\% 视为男胎检测达到基本可靠性的判据；
    \item 题目所给样本能够代表该地区高 BMI 孕妇群体的总体特征；
    \item 若某孕妇多次检测均未达到 4\%，则认为其在观测区间内未达标；
    \item 女胎异常判定中，附件 AB 列给出的非整倍体标记视为真实标签；
    \item 由于赛题未给出显式误差分布，本文用阈值邻域和分组达标比例变化来表征检测误差影响。
\end{enumerate}

\section{符号说明}
\begin{table}[H]
    \centering
    \begin{tabular}{>{\centering\arraybackslash}m{2.5cm} m{9.5cm}}
        \toprule
        符号 & 含义 \\
        \midrule
        $w$ & 检测孕周（以周为单位） \\
        $b$ & 孕妇 BMI \\
        $y$ & 男胎 Y 染色体浓度 \\
        $t_0$ & 某孕妇 Y 染色体浓度首次达到 4\% 的最早达标时间 \\
        $p_g(t)$ & 第 $g$ 组孕妇在孕周 $t$ 前达标的比例 \\
        $t_g^*$ & 第 $g$ 个 BMI 分组的最佳 NIPT 时点 \\
        $R(t)$ & 在时点 $t$ 下的潜在风险函数 \\
        $z_{13},z_{18},z_{21}$ & 13、18、21 号染色体的 Z 值 \\
        $x_c$ & 女胎 X 染色体浓度 \\
        $D$ & 女胎异常判定结果 \\
        \bottomrule
    \end{tabular}
\end{table}

\section{数据预处理}
\subsection{数据概况}
附件包含两张工作表：男胎检测数据与女胎检测数据。经统计：男胎数据共 1082 条记录，对应 267 名孕妇；女胎数据共 605 条记录，对应 147 名孕妇；男胎样本中胎儿健康 1044 条，不健康 38 条；女胎样本中带有非整倍体异常标签的记录共 67 条，正常记录 538 条。

\subsection{孕周标准化}
原始孕周以“11w+6”“20w+1”等形式给出。本文统一将其换算为十进制周数，例如“11w+6”记为 $11+6/7$ 周，便于后续回归和分组分析。

\subsection{重复检测处理}
针对男胎最早达标时间问题，本文按“孕妇代码”分组，并对同一孕妇的多条记录按孕周排序。若某次或某几次检测的 Y 染色体浓度达到 4\%，则取其最小孕周作为该孕妇的最早达标时间。该定义能有效抑制重复检测带来的偏差。

\section{问题一：Y 染色体浓度与孕周、BMI 等指标关系模型}
\subsection{相关性分析}
对男胎样本进行 Pearson 相关分析，得到 Y 染色体浓度与主要变量的相关系数如下：与孕周相关系数为 0.1266，与 BMI 相关系数为 -0.1513，与年龄相关系数为 -0.1194，与身高相关系数为 -0.1047，与体重相关系数为 -0.1806。结果表明，Y 染色体浓度随孕周增长而上升，随 BMI 和体重增加而下降。

\begin{figure}[H]
    \centering
    \includegraphics[width=0.72\textwidth]{figures/y_concentration_vs_week.pdf}
    \caption{男胎 Y 染色体浓度与孕周关系图}
    \label{fig:y-week}
\end{figure}

图 \ref{fig:y-week} 用于展示男胎 Y 染色体浓度随孕周增长的总体趋势，同时可直观表现高 BMI 样本的浓度偏低现象。

\subsection{两因素关系模型}
考虑题意强调 BMI 与孕周是主要因素，建立线性关系模型：
\[
y = \beta_0 + \beta_1 w + \beta_2 b + \varepsilon.
\]
最小二乘估计得到：
\[
\hat y = 0.11946 + 0.001255w - 0.001964b.
\]
模型拟合优度为 $R^2=0.0457$。其中，孕周项和 BMI 项均具有很高显著性，说明孕周每增加 1 周，Y 染色体浓度平均上升约 0.00126；BMI 每增加 1 个单位，Y 染色体浓度平均下降约 0.00196。

\subsection{多因素扩展模型}
进一步加入年龄、身高、体重后，建立五变量线性模型：
\[
y = \beta_0 + \beta_1 w + \beta_2 b + \beta_3 a + \beta_4 h + \beta_5 m + \varepsilon.
\]
计算得到模型 $R^2=0.0715$。其中孕周项显著正向，体重项与年龄项显著负向，身高项弱正向，BMI 项受共线性影响处于边界显著状态。这说明 BMI 适合用作主分组变量，而年龄、身高、体重可用于组内修正。

\subsection{问题一结论}
问题一表明：男胎 Y 染色体浓度与孕周显著正相关，与 BMI、体重、年龄呈负相关；以孕周和 BMI 为核心解释变量的关系模型能够较好反映总体趋势，并可作为后续时点优化的基础。

\section{问题二：基于 BMI 分组的最佳 NIPT 时点}
\subsection{最早达标时间定义}
定义某男胎孕妇的最早达标时间 $t_0$ 为其 Y 染色体浓度首次达到或超过 4\% 的最小孕周。267 名男胎孕妇中，大多数在观测区间内可达标，最早达标时间均值约为 13.64 周，中位数约为 12.86 周，并且随 BMI 升高而整体后移。

\subsection{分组原则}
比较多个候选分组方案后，本文采用如下四组：
\[
[0,30),\quad [30,34),\quad [34,38),\quad [38,+\infty).
\]
该方案兼顾了样本均衡性与分层差异性，各组样本数分别为 76、136、38 和 10。

\subsection{各组达标时间分布}
按上述分组统计最早达标时间，得到：
\begin{itemize}
    \item $[0,30)$ 组：平均 13.52 周，中位数 12.79 周；
    \item $[30,34)$ 组：平均 13.31 周，中位数 12.71 周；
    \item $[34,38)$ 组：平均 14.34 周，中位数 13.21 周；
    \item $[38,+\infty)$ 组：平均 15.50 周，中位数 15.86 周。
\end{itemize}

\begin{figure}[H]
    \centering
    \includegraphics[width=0.72\textwidth]{figures/bmi_group_reach_curve.pdf}
    \caption{不同 BMI 分组下的 Y 染色体浓度达标比例曲线}
    \label{fig:bmi-reach}
\end{figure}

图 \ref{fig:bmi-reach} 用于展示各 BMI 分组中“在给定孕周之前已达到 4\% 阈值”的比例变化，从而支持最佳检测时点选择。

\subsection{风险函数与最佳时点选择}
为平衡早检风险与延迟风险，构造简化风险函数：
\[
R_g(t)=\lambda_1\bigl(1-p_g(t)\bigr)+\lambda_2 t,
\]
其中 $p_g(t)$ 表示第 $g$ 组在孕周 $t$ 前达到 4\% 阈值的比例。本文采用“在较高达标率下尽量提前”的原则，即取达标比例接近或超过 80\% 的最早孕周作为推荐时点。由此得到：
\begin{itemize}
    \item BMI $[0,30)$：推荐检测时点 14 周；
    \item BMI $[30,34)$：推荐检测时点 14 周；
    \item BMI $[34,38)$：推荐检测时点 17 周；
    \item BMI $[38,+\infty)$：推荐检测时点 19\textasciitilde20 周，建议 19 周后尽快复核。
\end{itemize}

\subsection{检测误差影响分析}
检测误差可能使观测到的 Y 染色体浓度在阈值 4\% 附近波动，进而影响“最早达标时间”判定。对于低 BMI 组，由于多数样本在 14 周前后已稳定达标，阈值附近波动对推荐时点影响较小；对于高 BMI 组，部分样本长期处于阈值附近，因此微小误差可能导致达标时间推迟 1\textasciitilde2 周。可见 BMI 越高，检测误差对时点决策越敏感。

\section{问题三：综合多因素的最佳时点优化}
\subsection{模型思路}
问题三要求在 BMI 之外进一步考虑年龄、身高、体重等因素以及达标比例与检测误差。根据问题一回归结果，体重与年龄对 Y 染色体浓度具有显著负效应，因此本文保留 BMI 分组框架，同时引入多因素修正。设综合评分函数为：
\[
S = \alpha_1 b + \alpha_2 m + \alpha_3 a - \alpha_4 h.
\]
其中系数方向由回归符号确定：BMI、体重、年龄越高，达标风险越大；身高越高，风险相对减小。

\subsection{综合优化结果}
结合多因素模型与达标曲线，可给出如下推荐：
\begin{itemize}
    \item BMI $[0,30)$：基础时点 14 周；若年龄较大且体重偏高，则建议 15 周；
    \item BMI $[30,34)$：基础时点 14 周；高体重个体建议 15\textasciitilde16 周；
    \item BMI $[34,38)$：基础时点 17 周；若伴随高年龄、高体重，则建议 17\textasciitilde18 周；
    \item BMI $[38,+\infty)$：基础时点 19 周；若检测结果接近阈值，应进行 1 周内复测。
\end{itemize}

\subsection{检测误差影响}
在高 BMI 组中，14\textasciitilde18 周区间内达标比例增长较快，说明推荐时点对误差和个体差异较敏感。因此，问题三的关键并非简单追求更晚检测，而是通过多因素识别“易晚达标”的高风险人群，并对其采取更保守的时点安排。

\section{问题四：女胎异常判定方法}
\subsection{数据特点}
女胎数据共 605 条记录，其中 67 条带有 13、18、21 号染色体非整倍体异常标签，538 条为正常样本，异常比例约为 11.07\%。样本具有明显类别不平衡特征，因此判定模型应优先保证对异常样本的识别能力。

\subsection{特征比较}
对正常组与异常组的主要指标进行对比，发现：X 染色体浓度差异最为明显，正常组中位数约为 -0.0064，异常组中位数约为 -0.0223；单独考察 13、18、21 号染色体 Z 值时，异常组与正常组存在一定重叠，仅靠单阈值难以获得较高召回率；GC 含量、读段数及比对比例等质量指标对异常判定有辅助作用，但主导性弱于 X 染色体浓度。

\subsection{单一 Z 值规则的局限性}
设 $z_{\max}=\max(z_{13},z_{18},z_{21})$。当采用纯 Z 值阈值规则时，若阈值取 2.5，则召回率约为 11.94\%，准确率约为 79.67\%；若阈值取 3.0，则召回率仅约为 4.48\%，虽然准确率更高，但漏检过多。因此单纯依赖染色体 Z 值会显著降低异常样本识别能力。

\subsection{联合判定规则}
结合 X 染色体浓度分布差异，提出如下联合判定规则：
\[
D=1 \iff \bigl(x_c \le -0.015\bigr)\ \text{或}\ \max(z_{13},z_{18},z_{21})\ge 2.5.
\]
其中 $D=1$ 表示判定为异常。该规则的样本内效果为：真阳性 47，假阳性 193，真阴性 345，假阴性 20，召回率约 70.15\%，特异度约 64.13\%，总体准确率约 64.79\%。

\begin{figure}[H]
    \centering
    \includegraphics[width=0.72\textwidth]{figures/female_abnormal_rule.pdf}
    \caption{女胎异常联合判定规则示意图}
    \label{fig:female-rule}
\end{figure}

图 \ref{fig:female-rule} 用于展示 X 染色体浓度阈值和染色体 Z 值阈值共同作用下的判定逻辑，可作为论文中规则解释与结果展示的关键图形。

\subsection{方法解释}
该规则的临床解释为：当 X 染色体浓度明显偏低时，说明女胎相关信号存在异常偏移；当 13、18、21 号染色体中任一 Z 值较高时，对应非整倍体风险上升；将两者联合，可弥补单一 Z 值模型召回率不足的缺陷。

\section{模型评价}
\subsection{模型优点}
\begin{enumerate}
    \item 可解释性强，各问模型均围绕孕周、BMI、Y 浓度阈值和染色体指标展开；
    \item 与题意贴合，第二问和第三问均以“最早达标时间”和“潜在风险最小”为核心；
    \item 处理了重复检测问题，通过按孕妇聚合估计最早达标时间，避免了重复样本干扰；
    \item 对女胎异常判定具有可操作性，联合阈值规则简单清晰。
\end{enumerate}

\subsection{模型不足}
\begin{enumerate}
    \item 线性模型对复杂非线性关系的刻画能力有限，拟合优度不高；
    \item 高 BMI 组样本量较少，导致分组边界和最佳时点的稳定性受限；
    \item 女胎异常判定使用的是规则型模型，尚未结合更复杂的概率模型或交叉验证策略；
    \item 未获得真实临床误差分布，只能通过阈值邻域分析误差影响。
\end{enumerate}

\section{模型改进与推广}
\begin{enumerate}
    \item 对问题一可进一步建立广义加性模型或混合效应模型，以处理同一孕妇多次检测的相关性；
    \item 对问题二和问题三可引入生存分析或分位数回归，以更精确描述“首次达标时间”；
    \item 对问题四可在更多样本下构造代价敏感分类模型，以提高异常样本识别率；
    \item 若未来获得更详细的临床标注，还可将本模型推广到更细粒度的个体化 NIPT 时点推荐系统。
\end{enumerate}

\section{结论}
本文围绕 NIPT 的时点选择与胎儿异常判定问题，构建了较完整的数学建模框架。研究表明：男胎 Y 染色体浓度随孕周增加而升高，随 BMI 升高而降低；基于 BMI 的四组划分 $[0,30)$、$[30,34)$、$[34,38)$、$[38,+\infty)$ 能较好刻画不同孕妇群体的达标时间差异，相应推荐检测时点分别为 14 周、14 周、17 周和 19 周左右。进一步结合年龄、身高、体重等因素可对推荐时点进行组内修正，提高个体化决策能力。对于女胎异常判定，单纯依赖染色体 Z 值效果有限，而采用 X 染色体浓度与染色体 Z 值联合规则，可显著提升异常样本识别能力。整体而言，本文模型为高 BMI 孕妇群体的 NIPT 时点制定与女胎异常初筛提供了量化支持。

\begin{thebibliography}{99}
\bibitem{ref1} 全国大学生数学建模竞赛组委会. 2025 年高教社杯全国大学生数学建模竞赛 C 题.
\bibitem{ref2} 全国大学生数学建模竞赛组委会. 全国大学生数学建模竞赛论文格式规范（2026 年修订稿）.
\bibitem{ref3} 无创产前检测相关临床文献与公开资料.
\end{thebibliography}

\appendix
\section{主要统计结果}
男胎数据：1082 条记录，267 名孕妇；女胎数据：605 条记录，147 名孕妇。

两因素线性模型：
\[
\hat y = 0.11946 + 0.001255w - 0.001964b.
\]

推荐 BMI 分组与检测时点为：$[0,30)$ 取 14 周，$[30,34)$ 取 14 周，$[34,38)$ 取 17 周，$[38,+\infty)$ 取 19 周。女胎联合判定规则为：$x_c\le -0.015$ 或 $\max(z_{13},z_{18},z_{21})\ge 2.5$。

\end{document}

